\documentclass[11pt]{article}

\usepackage[margin=1in]{geometry}
\usepackage{amsmath,amssymb,amsthm,mathtools}
\usepackage{booktabs,tabularx}
\usepackage[hidelinks]{hyperref}

\newtheorem{theorem}{Theorem}
\newtheorem{proposition}[theorem]{Proposition}
\newtheorem{corollary}[theorem]{Corollary}
\newtheorem{lemma}[theorem]{Lemma}
\newtheorem{counterexample}[theorem]{Counterexample}
\theoremstyle{definition}
\newtheorem{definition}[theorem]{Definition}
\theoremstyle{remark}
\newtheorem{remark}[theorem]{Remark}

\newcommand{\R}{\mathbb{R}}
\newcommand{\E}{\mathbb{E}}
\newcommand{\Prob}{\mathbb{P}}
\newcommand{\KL}{\operatorname{KL}}
\newcommand{\Cat}{\operatorname{Categorical}}
\newcommand{\1}{\mathbf{1}}
\newcommand{\dd}{\,\mathrm{d}}
\newcommand{\abs}[1]{\left|#1\right|}
\newcommand{\norm}[1]{\left\lVert#1\right\rVert}
\setlength{\emergencystretch}{3em}

\title{Tempered Ensembles of Reverse-KL Neural Transports\\
for Multimodal Hamiltonian Monte Carlo}
\author{BayesFilter research record}
\date{Revised 28 August 2026}

\begin{document}
\maketitle

\begin{abstract}
This note corrects the proposed role of adaptive particle replay in training a
NeuTra map for multimodal SSL--LSTM posteriors.  Proposal-aware replay can be
mathematically consistent under strong support, moment, and stochastic-
approximation assumptions, but it does not solve the finite high-dimensional
discovery problem.  If a learned transport generates the target-side measure
used to train and validate itself, an operationally missing mode remains
unobserved and the argument becomes circular.  We retain the original NeuTra
reverse Kullback--Leibler objective with fresh independent Gaussian draws and
replace the single learned map by a probability mixture of independently
initialized transport maps trained along a proper temperature path.  We prove
the mixture-density and reverse-KL identities, characterize when variational
mixture weights equal mode masses, quantify their approximation-error bias,
and prove invariance of a frozen multi-chart HMC kernel embedded in replica
exchange.  The construction avoids a particle approximation as a foundation,
but it does not guarantee discovery, irreducibility, mixing, or favorable
high-dimensional scaling.
\end{abstract}

\section{Status, correction, and claim boundary}

This LaTeX document is the formal audit surface for the 2026-08-28 revision of
\path{docs/plans/bayesfilter-ssl-lstm-q20-adaptive-replay-neutra-mathematical-note-2026-08-21.md}.
The replay identities and conditional convergence theorems in
the original note are not retracted.  The retracted proposition is the
\emph{scientific recommendation} that an adaptively generated particle measure
should be the primary high-dimensional repair for NeuTra.

The proposed replacement has two layers.  During training, all stochastic
queries remain draws from a standard Gaussian and optimize a reverse-KL
objective.  During posterior sampling, every learned map is frozen and used
only as a coordinate chart for an exact Metropolis-corrected HMC kernel.
Replica exchange supplies the tempered transition mechanism.  The final cold
marginal is exact under the assumptions proved below; its finite run need not
have mixed.

The current q=20 target has four physical parameters even though its likelihood
is evaluated by a higher-dimensional filtering computation.  This note is
motivated by future SSL--LSTM targets with many more parameters.  It contains
no claim that the proposed method avoids the curse of dimensionality or that a
finite number of charts discovers every important mode.

\section{Target, transports, and reverse KL}

Let $\Theta=\R^d$.  Let $\widetilde\pi:\Theta\to(0,\infty)$ be measurable with
$0<Z=\int\widetilde\pi(\theta)\dd\theta<\infty$, and write
$\pi=\widetilde\pi/Z$.  Let $\rho$ be the standard Gaussian density on
$\R^d$.  For each component $i\in\{1,\ldots,K\}$, let
$T_i:\R^d\to\R^d$ be a $C^1$ diffeomorphism with nonzero Jacobian determinant.
Its pushforward density is
\begin{equation}
 q_i(\theta)
 =\rho\{T_i^{-1}(\theta)\}
  \abs{\det D T_i^{-1}(\theta)}.
 \label{eq:component-density}
\end{equation}

\begin{proposition}[Original NeuTra reverse-KL identity]
\label{prop:reverse-kl}
For one transport $T_i$,
\begin{equation}
 \KL(q_i\Vert\pi)
 =\E_{Z_0\sim\rho}\!\left[
   \log\rho(Z_0)-\log\abs{\det DT_i(Z_0)}
   -\log\widetilde\pi\{T_i(Z_0)\}
 \right]+\log Z.
 \label{eq:reverse-kl}
\end{equation}
Thus its trainable term requires independent Gaussian base draws and target
evaluations, but no draws from $\pi$.
\end{proposition}

\begin{proof}
Set $\theta=T_i(z)$ in the defining integral for
$\KL(q_i\Vert\pi)$.  The change-of-variables identity gives
$q_i\{T_i(z)\}\abs{\det DT_i(z)}=\rho(z)$.  Substitution and
$\log\pi=\log\widetilde\pi-\log Z$ yield
\eqref{eq:reverse-kl}.
\end{proof}

Under a local integrable-envelope condition, differentiation can be moved
through the Gaussian expectation.  A fresh Gaussian minibatch then gives an
unbiased stochastic gradient of the displayed trainable objective.  This is
the paper-style NeuTra objective, not the particle-trained forward-KL
objective in the current weighted trainer.

\section{Why proposal support is insufficient}

Let $m_d$ be a normalized proposal for a normalized target $\pi_d$, and set
$w_d=\pi_d/m_d$.  If $\pi_d\ll m_d$, then
\begin{equation}
 \E_{m_d}w_d=1,
 \qquad
 \E_{m_d}w_d^2
 =1+\chi^2(\pi_d\Vert m_d)
 =\exp\{D_2(\pi_d\Vert m_d)\}.
 \label{eq:weight-second-moment}
\end{equation}
The reciprocal of the second moment is a population relative-ESS proxy.  It is
not an exact finite-sample ESS identity.

\begin{proposition}[Product mismatch]
\label{prop:product-mismatch}
Suppose
\[
 \pi_d(\theta)=\prod_{j=1}^d\pi_j(\theta_j),
 \qquad
 m_d(\theta)=\prod_{j=1}^d m_j(\theta_j),
\]
with $\pi_j\ll m_j$, and define
$c_j=\int\pi_j(x)^2/m_j(x)\dd x$.  Then
\begin{equation}
 \E_{m_d}w_d^2=\prod_{j=1}^d c_j.
 \label{eq:product-second-moment}
\end{equation}
If $c_j\geq1+\delta$ for all $j$ and some $\delta>0$, then
\begin{equation}
 \{\E_{m_d}w_d^2\}^{-1}\leq(1+\delta)^{-d}.
 \label{eq:exponential-ess}
\end{equation}
\end{proposition}

\begin{proof}
The nonnegative integrand factorizes:
\[
 w_d(\theta)^2m_d(\theta)
 =\prod_{j=1}^d\frac{\pi_j(\theta_j)^2}{m_j(\theta_j)}.
\]
Tonelli's theorem yields \eqref{eq:product-second-moment}.  Bounding every
factor below by $1+\delta$ and taking the reciprocal proves
\eqref{eq:exponential-ess}.
\end{proof}

\begin{proposition}[No correction without a sampled row]
\label{prop:no-row-no-correction}
For a measurable $A$ and independent $X_1,\ldots,X_N\sim m$,
\begin{equation}
 \Prob\{X_i\notin A\text{ for all }i\}=\bigl\{1-m(A)\bigr\}^N.
 \label{eq:missing-region}
\end{equation}
On this event, every weighted atomic measure supported on the sampled rows
assigns zero mass to $A$.
\end{proposition}

\begin{proof}
Independence proves \eqref{eq:missing-region}.  The indicator of $A$ is zero at
every support point of the empirical measure, regardless of its weights.
\end{proof}

These propositions identify the circularity.  A defensive proposal component
can establish $m(A)>0$ without making $Nm(A)$ appreciable or the second moment
in \eqref{eq:weight-second-moment} finite.  If $m$ is primarily constructed
from a candidate $q_\phi$, the same candidate cannot obtain independent global
validation from a finite measure generated under $m$.

\section{Probability mixtures of transports}

Let $\alpha=(\alpha_1,\ldots,\alpha_K)$ lie in the positive probability
simplex.

\begin{proposition}[Mixture pushforward]
\label{prop:mixture-pushforward}
Draw $I\sim\Cat(\alpha)$, independently draw $Z_0\sim\rho$, and set
$\Theta=T_I(Z_0)$.  Then
\begin{equation}
 q_\alpha(\theta)=\sum_{i=1}^K\alpha_i q_i(\theta).
 \label{eq:mixture-density}
\end{equation}
\end{proposition}

\begin{proof}
For any measurable $A$, condition on $I$:
\[
 \Prob(\Theta\in A)
 =\sum_i\alpha_i\Prob\{T_i(Z_0)\in A\}
 =\sum_i\alpha_i\int_Aq_i(\theta)\dd\theta.
\]
This identifies \eqref{eq:mixture-density}.
\end{proof}

\begin{counterexample}[A weighted average of maps is not the mixture]
\label{counter:averaged-map}
In one dimension, let $T_1(z)=z-a$, $T_2(z)=z+a$, and
$\alpha_1=\alpha_2=1/2$.  The averaged map is $\overline T(z)=z$ and pushes a
standard normal to $N(0,1)$.  Equation~\eqref{eq:mixture-density} instead gives
$\frac12N(-a,1)+\frac12N(a,1)$, whose variance is $1+a^2$.  These laws differ
for $a\neq0$.  Averaging diffeomorphisms can also destroy invertibility:
the average of $z$ and $-z$ is the constant zero map.
\end{counterexample}

The ensemble therefore has a discrete chart index.  It is not one invertible
NeuTra map.

\begin{proposition}[Mixture reverse KL from Gaussian draws]
\label{prop:mixture-reverse-kl}
For a proper target $\pi_\beta=\widetilde\pi_\beta/Z_\beta$,
\begin{align}
 \KL(q_\alpha\Vert\pi_\beta)
 &=\sum_{i=1}^K\alpha_i\E_{Z_0\sim\rho}\left[
    \log q_\alpha\{T_i(Z_0)\}
    -\log\widetilde\pi_\beta\{T_i(Z_0)\}
   \right]+\log Z_\beta.
 \label{eq:mixture-reverse-kl}
\end{align}
No target-distributed particle is required to evaluate the trainable term.
\end{proposition}

\begin{proof}
Split the KL expectation according to
$q_\alpha=\sum_i\alpha_iq_i$ and then apply the change of variables
$\theta=T_i(z)$ to each component expectation.  Replacing
$\log\pi_\beta$ by $\log\widetilde\pi_\beta-\log Z_\beta$ proves the result.
\end{proof}

When the finite component sum is enumerated and the required derivatives have
integrable envelopes, fresh Gaussian batches give a stochastic gradient of
\eqref{eq:mixture-reverse-kl}.  Evaluation costs scale at least as the required
cross-component density calculations; the identity is not a computational
scaling theorem.

\section{Variational weights and posterior mode masses}

Let $A_1,\ldots,A_K$ be a measurable partition with
$p_i=\pi(A_i)>0$, and set
$\pi_i=\pi\1_{A_i}/p_i$.  Suppose $q_i$ is supported in $A_i$ and define
$\delta_i=\KL(q_i\Vert\pi_i)$.  This is an idealized separated-region model;
full-support neural flows approximate rather than literally satisfy it.

\begin{proposition}[Separated-region decomposition]
\label{prop:separated-weights}
Under the preceding assumptions,
\begin{equation}
 \KL(q_\alpha\Vert\pi)
 =\sum_{i=1}^K\alpha_i
   \left\{\log\frac{\alpha_i}{p_i}+\delta_i\right\}.
 \label{eq:separated-kl}
\end{equation}
For fixed components, the unique minimizing weights are
\begin{equation}
 \alpha_i^*
 =\frac{p_i e^{-\delta_i}}
        {\sum_{j=1}^Kp_j e^{-\delta_j}}.
 \label{eq:biased-mixture-weight}
\end{equation}
\end{proposition}

\begin{proof}
On $A_i$, $q_\alpha=\alpha_iq_i$ and $\pi=p_i\pi_i$.  Splitting the KL
integral over the partition proves \eqref{eq:separated-kl}.  Differentiating
with a multiplier for $\sum_i\alpha_i=1$ gives
\[
 \log(\alpha_i/p_i)+\delta_i+1+\lambda=0.
\]
Normalization yields \eqref{eq:biased-mixture-weight}; strict convexity on the
positive simplex gives uniqueness.
\end{proof}

\begin{corollary}[Exact components recover regional masses]
\label{cor:exact-regional-masses}
If $q_i=\pi_i$ for all $i$, then $\delta_i=0$ and
$\alpha_i^*=p_i$.  If the $\delta_i$ differ, the optimizing variational weights
confound regional mass and local approximation error.
\end{corollary}

This distinction is essential.  The $\alpha_i$ are mixture-density parameters.
They are not posterior mode-mass authority unless the regional approximation
conditions are established.

\section{Limits of random-start discovery}

\begin{proposition}[Conditional multi-start coverage]
\label{prop:multistart}
Suppose one independently initialized training run reaches the attraction basin
associated with declared mode $j$ with probability $r_j$.  With $K$
independent runs, the probability of missing mode $j$ is
\begin{equation}
 (1-r_j)^K.
 \label{eq:mode-miss-probability}
\end{equation}
For a finite declared collection of $J$ modes, the probability of missing at
least one is no greater than
$\sum_{j=1}^J(1-r_j)^K$.
\end{proposition}

\begin{proof}
The first equality follows from independence.  The second statement is the
union bound over the declared modes.
\end{proof}

The result is conditional, not a coverage theorem: $r_j$ is unknown and can be
zero.  Different neural-network seeds do not necessarily produce different
physical initial locations.

\begin{proposition}[Finite-query non-identification]
\label{prop:finite-query}
Let an algorithm query a smooth positive integrable unnormalized target and
finitely many of its derivatives at a finite set $S\subset\R^d$.  There is
another smooth positive integrable target that agrees with all queried values
and derivatives but places arbitrarily large additional unnormalized mass in a
region disjoint from $S$.
\end{proposition}

\begin{proof}
Choose an open ball $B$ whose closure misses $S$, and choose a nonzero,
nonnegative $C^\infty$ bump $h$ supported in $B$.  For the original target
$\widetilde\pi_0$, set
\begin{equation}
 \widetilde\pi_c=\widetilde\pi_0+c h,
 \qquad c>0.
 \label{eq:hidden-bump}
\end{equation}
The bump and all derivatives vanish near every point in $S$, while its added
mass is $c\int h$, which is arbitrarily large as $c$ increases.  Positivity,
smoothness, and integrability are preserved.
\end{proof}

No finite generic computation can certify exhaustive global discovery without
additional structure.  This limitation applies to transport ensembles,
tempering, MCMC, and particle methods alike.

\section{Proper tempering and reverse-KL continuation}

Let $g_0$ be a normalized positive reference.  Define
\begin{equation}
 \widetilde\pi_\beta(\theta)
 =g_0(\theta)^{1-\beta}\widetilde\pi(\theta)^\beta,
 \qquad 0\leq\beta\leq1,
 \label{eq:geometric-bridge}
\end{equation}
and assume $0<Z_\beta<\infty$ throughout.  If
$\widetilde\pi(\theta)=p(\theta)L(y\mid\theta)$ with proper prior $p$, choosing
$g_0=p$ yields
\begin{equation}
 \widetilde\pi_\beta(\theta)=p(\theta)L(y\mid\theta)^\beta.
 \label{eq:likelihood-tempering}
\end{equation}
An unnormalized uniform endpoint on $\R^d$ is not an admissible $\beta=0$
reference.

\begin{proposition}[Energy interpolation]
\label{prop:energy-interpolation}
Let $U_0=-\log g_0$ and $U_1=-\log\widetilde\pi$.  Up to an additive constant,
\begin{equation}
 U_\beta(\theta)=(1-\beta)U_0(\theta)+\beta U_1(\theta).
 \label{eq:energy-interpolation}
\end{equation}
For any $a,b\in\R^d$,
\begin{align}
 U_\beta(b)-U_\beta(a)
 &=(1-\beta)\{U_0(b)-U_0(a)\}
   +\beta\{U_1(b)-U_1(a)\}.
 \label{eq:barrier-interpolation}
\end{align}
For likelihood tempering, only the log-likelihood contribution is multiplied
by $\beta$.
\end{proposition}

\begin{proof}
Take negative logarithms in \eqref{eq:geometric-bridge} and subtract the values
at $a$ and $b$.  The factorization in \eqref{eq:likelihood-tempering} gives the
last statement.
\end{proof}

Equation~\eqref{eq:barrier-interpolation} explains why tempering can reduce a
likelihood barrier.  It also rules out a generic monotonic-connectivity claim:
the reference energy and regional volumes can change the relevant geometry and
mass ordering.

\begin{proposition}[The cold optimization target is unchanged]
\label{prop:cold-objective}
If the final inverse temperature is $\beta_L=1$, then the final mixture
objective is exactly $\KL(q_\alpha\Vert\pi)$.  Earlier temperatures and warm
starts can change the optimization path but not this population objective.
\end{proposition}

\begin{proof}
Equation~\eqref{eq:geometric-bridge} gives
$\widetilde\pi_1=\widetilde\pi$, hence $\pi_1=\pi$.  Substitute into
Proposition~\ref{prop:mixture-reverse-kl}.
\end{proof}

Tempering is therefore an optimization and discovery mechanism during training,
not evidence that the optimizer found distinct modes.

\section{Frozen transports as exact HMC coordinate charts}

Fix $\beta$ and chart $T_i$.  Define the exact pullback density
\begin{equation}
 \pi_{\beta i}^z(z)
 =\pi_\beta\{T_i(z)\}\abs{\det DT_i(z)}.
 \label{eq:chart-pullback}
\end{equation}
Let $P_{\beta i}$ be a Markov kernel preserving
$\pi_{\beta i}^z$, such as a fixed Metropolis-corrected HMC kernel using the
exact value and score of \eqref{eq:chart-pullback}.  Define
\begin{equation}
 K_{\beta i}(\theta,A)
 =P_{\beta i}\{T_i^{-1}(\theta),T_i^{-1}(A)\}.
 \label{eq:physical-chart-kernel}
\end{equation}

\begin{proposition}[Chart-kernel invariance]
\label{prop:chart-invariance}
The physical kernel $K_{\beta i}$ preserves $\pi_\beta$.
\end{proposition}

\begin{proof}
The change of variables $\theta=T_i(z)$ sends $\pi_\beta$ to
$\pi_{\beta i}^z$.  Hence
\begin{align*}
 \int\pi_\beta(\dd\theta)K_{\beta i}(\theta,A)
 &=\int\pi_{\beta i}^z(\dd z)
      P_{\beta i}\{z,T_i^{-1}(A)\}\\
 &=\pi_{\beta i}^z\{T_i^{-1}(A)\}
 =\pi_\beta(A).
\end{align*}
\end{proof}

\begin{proposition}[Fixed mixtures of invariant chart kernels]
\label{prop:kernel-mixture}
If $\gamma_i\geq0$, $\sum_i\gamma_i=1$, and the $\gamma_i$ are fixed and
state-independent, then
\begin{equation}
 K_\beta=\sum_i\gamma_iK_{\beta i}
 \label{eq:kernel-mixture}
\end{equation}
preserves $\pi_\beta$.
\end{proposition}

\begin{proof}
By linearity and Proposition~\ref{prop:chart-invariance},
$\pi_\beta K_\beta=\sum_i\gamma_i\pi_\beta=\pi_\beta$.
\end{proof}

\begin{counterexample}[State-dependent chart selection]
\label{counter:state-dependent-mixture}
On the uniform target over $\{0,1\}$, both the identity kernel and the flip
kernel preserve the target.  Select the identity at state zero and the flip at
state one.  Both states then transition to zero, so the state-dependent mixture
does not preserve the uniform law.
\end{counterexample}

State-dependent chart selection requires a separate Metropolis, Gibbs, or
augmented-state argument.  It is excluded from the initial proposal.

\section{Replica exchange with chart ensembles}

Let $0\leq\beta_0<\cdots<\beta_L=1$ and define
\begin{equation}
 \Pi(\theta_0,\ldots,\theta_L)
 =\prod_{\ell=0}^L\pi_{\beta_\ell}(\theta_\ell).
 \label{eq:product-target}
\end{equation}
At level $\ell$, let $K_\ell$ be a fixed mixture of exact chart kernels from
Proposition~\ref{prop:kernel-mixture}.

\begin{proposition}[Adjacent replica exchange]
\label{prop:replica-swap}
A symmetric proposal to exchange $\theta_\ell$ and $\theta_{\ell+1}$, accepted
with probability
\begin{equation}
 a_{\mathrm{swap}}
 =1\wedge
 \frac{
  \widetilde\pi_{\beta_\ell}(\theta_{\ell+1})
  \widetilde\pi_{\beta_{\ell+1}}(\theta_\ell)}{
  \widetilde\pi_{\beta_\ell}(\theta_\ell)
  \widetilde\pi_{\beta_{\ell+1}}(\theta_{\ell+1})},
 \label{eq:swap-ratio}
\end{equation}
is reversible with respect to $\Pi$.
\end{proposition}

\begin{proof}
The ratio of \eqref{eq:product-target} at the exchanged and current states
cancels every unaffected factor and both affected normalizing constants.  The
remainder is \eqref{eq:swap-ratio}, the ordinary Metropolis ratio for a
symmetric involution.
\end{proof}

\begin{theorem}[Exact cold marginal]
\label{thm:exact-cold-marginal}
Suppose every chart is frozen, every within-chart HMC kernel preserves the
exact pullback \eqref{eq:chart-pullback}, every chart-selection mixture is fixed
and state-independent, and every adjacent exchange uses
\eqref{eq:swap-ratio}.  Any fixed composition or state-independent random scan
of the within-temperature and exchange kernels preserves $\Pi$.  Its
$\beta_L=1$ marginal is $\pi$.
\end{theorem}

\begin{proof}
The product of the within-temperature kernels preserves
\eqref{eq:product-target}, since each factor kernel preserves its corresponding
factor.  Proposition~\ref{prop:replica-swap} gives the same invariance for each
swap.  Compositions and fixed state-independent mixtures of kernels with a
common invariant law preserve that law.  The final factor is
$\pi_{\beta_L}=\pi$.
\end{proof}

\begin{proposition}[Exactness is not discovery]
\label{prop:exact-not-discovery}
The identity kernel preserves every probability law but never changes state.
More generally, if every within-temperature kernel leaves a common region $A$
closed and all replicas begin in $A$, replica exchange cannot move a replica to
$A^c$.
\end{proposition}

\begin{proof}
The first statement is immediate.  Under the second hypothesis, local kernels
retain every state in $A$, while swaps only permute the existing states.
Induction over transitions proves the claim.
\end{proof}

Thus finite values, local acceptance, or adjacent swaps do not establish global
mixing.  Round trips through the temperature ladder, hot-level basin forgetting, repeated
cold-level transitions, initialization forgetting, and ordinary retained-chain
diagnostics are separate evidence requirements.

\section{Optional independence proposals}

\begin{proposition}[Metropolis-corrected mixture proposal]
\label{prop:mixture-independence-mh}
Assume $q_\alpha>0$ wherever $\pi>0$.  Propose $Y\sim q_\alpha$ independently
of current state $X=x$ and accept with
\begin{equation}
 a_{\mathrm{ind}}(x,y)
 =1\wedge
 \frac{\widetilde\pi(y)q_\alpha(x)}
      {\widetilde\pi(x)q_\alpha(y)}.
 \label{eq:independence-ratio}
\end{equation}
The resulting kernel is reversible with respect to $\pi$.
\end{proposition}

\begin{proof}
For $x\neq y$, the accepted probability flow is
\[
 \pi(x)q_\alpha(y)a_{\mathrm{ind}}(x,y)
 =\min\{\pi(x)q_\alpha(y),\pi(y)q_\alpha(x)\},
\]
which is symmetric in $x$ and $y$.  Rejection completes detailed balance.
\end{proof}

The correction establishes the invariant target, not useful acceptance.  In
high dimension this kernel can fail for the same mismatch quantified in
Proposition~\ref{prop:product-mismatch}.  It is therefore an optional arm, not
the primary sampler.

\section{Proposed algorithms}

\subsection{Tempered ensemble reverse-KL training}

\begin{enumerate}
 \item Expose and verify a proper reference-target bridge, preferably the
 prior-likelihood path \eqref{eq:likelihood-tempering}.
 \item Predeclare a bounded pilot temperature ladder.  Its number and spacing
 are hypotheses rather than defaults.
 \item Initialize several transports with independent stateless neural seeds
 and deliberately distinct affine locations and scales under the reference
 law.  Neural weight seeds alone do not establish physical diversity.
 \item At each temperature, train every lineage using the original Gaussian
 reverse-KL objective.  Warm-start from the preceding level and use only
 predeclared perturbation or branching rules.  Preserve every lineage through
 calibration rather than pruning solely by descriptive loss.
 \item Optionally refine all components and positive mixture weights jointly
 using \eqref{eq:mixture-reverse-kl}.  Treat the weights according to
 Proposition~\ref{prop:separated-weights}.
 \item Freeze all transports, mixture weights, temperature identities, target
 signatures, and selection rules before sampler validation.
\end{enumerate}

Full optimization at $\beta=0$ makes every component distribution target the
same reference law and can erase distributional diversity.  The endpoint is
therefore a bridge and implementation check, not the sole source of distinct
lineages.  The bounded design must include predeclared fresh restarts or
branching at one or more positive temperatures and must compare them with pure
continuation.  This changes initialization and optimization paths, not the
reverse-KL objective.

\subsection{Exact tempered multi-chart HMC}

\begin{enumerate}
 \item Construct the exact pullback \eqref{eq:chart-pullback} for every retained
 temperature/chart pair, including value, score, and Jacobian terms.
 \item Tune a fixed Metropolis-corrected HMC kernel in each declared scope using
 disjoint tuning draws.  Do not retain warmup as posterior output.
 \item Select charts with fixed state-independent frequencies, apply the chart
 inverse to the current physical state, run one exact transformed HMC update,
 and map the result back.
 \item Apply alternating adjacent swaps using \eqref{eq:swap-ratio}, recording
 replica identity and complete swap telemetry.
 \item Use only retained $\beta=1$ states for posterior summaries.  Require both
 ordinary sequential-HMC diagnostics and tempering-specific travel evidence.
 \item Evaluate Proposition~\ref{prop:mixture-independence-mh} only as a separate
 optional global-move arm.
\end{enumerate}

The current public BayesFilter fixed-transport tuner binds one frozen transport
and one transformed target per tuning artifact.  Running that tuner separately
for each $(\beta,i)$ pair would produce component artifacts, but it would not by
itself create the multi-chart sequential controller required above.  That
controller, its route-ledger entry, and its replay of the exact per-scope tuning
artifacts are new implementation work.

\section{Evidence roles and nonconclusions}

Three kinds of weights must not be conflated:
\begin{center}
\begin{tabularx}{\textwidth}{@{}lX@{}}
\toprule
Symbol & Meaning and claim boundary \\
\midrule
$\alpha_i$ & Mixture-density parameter in $q_\alpha$; equals a regional target
mass only under the separated exact-component conditions. \\
$\gamma_{\ell i}$ & Fixed frequency for selecting a chart kernel at temperature
$\ell$; affects efficiency but not the invariant posterior. \\
$p_i$ & Actual posterior probability of a declared region; must be inferred
from valid stationary posterior evidence, not assigned from initialization
counts. \\
\bottomrule
\end{tabularx}
\end{center}

The required comparison ladder is: a single cold reverse-KL transport; physical-
coordinate replica exchange under the same bridge; a cold multi-start ensemble;
a tempered ensemble without joint mixture refinement; and the jointly refined
tempered ensemble.  Reverse-KL loss, latent moments, component separation,
acceptance, and adjacent swap rates are explanatory.  Exact density/Jacobian/
score identities are implementation vetoes.  Replica round trips, hot-level
basin forgetting, cold retained rank-normalized and folded $\widehat R$, bulk
and tail ESS, direct mode transitions, and downstream reference agreement are
sampler gates under a predeclared experiment contract.

No theorem here establishes exhaustive mode discovery, geometric ergodicity,
favorable scaling in $d$, a bounded number of temperatures or components,
statistical superiority, exact inference for the nonlinear state-space model,
or correctness of the current finite SSL--LSTM implementation.  A future
dimension ladder must measure how components, temperatures, target calls,
memory, and wall time scale.

\section{Source and implementation boundary}

Proposition~\ref{prop:reverse-kl} follows the change-of-variables and reverse-
KL construction in Section~2.2, equations~(2)--(3), and Section~2.3 of Hoffman
et al.~\cite{hoffman2019}.  The product target and swap ratio in
Proposition~\ref{prop:replica-swap} are the construction in Section~II,
equations~(2.1)--(2.7), of Hukushima and Nemoto~\cite{hukushima1996}.  The use
of a frozen map as an exactly corrected MCMC coordinate transformation follows
Section~3.1, especially equation~(21), of Parno and Marzouk~\cite{parno2018}
and Section~2.3 of Hoffman et al.  The mixture identities, separated-region
weight decomposition, finite-query result, and multi-chart kernel composition
are project derivations.

The local inspected sources are recorded in
\path{.localresources/papers/multimodal_hmc/CORPUS_AUDIT.md}.  Relevant
implementation anchors are
\path{bayesfilter/inference/neutra_weighted_training.py} for the existing
transport and matched reverse-KL comparator, and
\path{bayesfilter/nonlinear/ssl_lstm_complexity_batched_target_tf.py} for the
q=20 batch-native likelihood-plus-Gaussian-prior target.  The active single-map
HMC constraints are documented in
\path{docs/reference/hmc-tuning-interface.md} and implemented under
\path{bayesfilter/inference/neutra_hmc.py} and the fixed-transport tuner.
The file \path{bayesfilter/testing/distributed_replica_exchange_tf.py} is a
diagnostic physical-coordinate replica-exchange implementation.  The latter is
not a production multi-chart implementation and its pure power-tempering path
does not by itself satisfy the proper-reference bridge required here.

\begin{thebibliography}{9}

\bibitem{hoffman2019}
M. D. Hoffman, P. Sountsov, J. V. Dillon, I. Langmore, D. Tran, and
S. Vasudevan.
\newblock NeuTra-lizing bad geometry in Hamiltonian Monte Carlo using neural
transport.
\newblock \emph{Proceedings of Machine Learning Research}, 96:1--17, 2019.

\bibitem{hukushima1996}
K. Hukushima and K. Nemoto.
\newblock Exchange Monte Carlo method and application to spin glass
simulations.
\newblock \emph{Journal of the Physical Society of Japan}, 65(6):1604--1608,
1996.

\bibitem{parno2018}
M. D. Parno and Y. M. Marzouk.
\newblock Transport map accelerated Markov chain Monte Carlo.
\newblock \emph{SIAM/ASA Journal on Uncertainty Quantification}, 6(2):645--682,
2018.

\end{thebibliography}

\end{document}
